% ===== Supplement (DRAFT, not \input yet): a theory of mode selection for hierarchical heads on frozen features =====
% Numbers: outputs/theory/{theory_checks.json, exp132_prediction_*.json}; scripts/theory_checks.py, scripts/theory_depth_prediction.py.
% Cell radii: exact, over all cell corners (script inline in the 2026-11-20 session; reproduced by tfpose.tesseract.cell_corners).
% Tags: [proved] exact statement with proof; [measured] number from our runs; [predicted] stated before the data exist; [heuristic].
\makeatletter
\@ifundefined{theorem}{\newtheorem{theorem}{Theorem}}{}
\@ifundefined{proposition}{\newtheorem{proposition}{Proposition}}{}
\@ifundefined{corollary}{\newtheorem{corollary}{Corollary}}{}
\makeatother

\section{Why Selection Dominates: A Theory of the Design}
\label{sec:supp_theory}

The title of the paper is a claim: with frozen features and small heads, the error of a rotation estimate is mostly the error
of \emph{selecting a mode}, not of locating the pose within the right mode. This section makes that claim precise and derives
the design from it. \Cref{thm:decomp}, the central result, splits the error of a tesseract tree into a \emph{selection}
term and a \emph{resolution} term and gives the exact effect of the tree depth on each. The propositions then place every
component of the method on one of the two terms: the hierarchy decides selection in its first levels
(\cref{prop:toplevels}), multi-path targets are the only targets of the three we use that are both symmetry-invariant and
consistent with the roll (\cref{prop:multipath}), orbit voting changes the information the frozen backbone delivers and has
a measurable failure condition (\cref{prop:orbit}), and the allocentric conversion is insensitive to translation error
(\cref{prop:allo}). \Cref{sec:supp_theory_scaling} turns \cref{thm:decomp} into a scaling law in the depth, which we state as a
prediction for the depth sweep (EXP-132) before any depth-6 or depth-7 tree was evaluated. Each statement is tagged
\emph{proved}, \emph{measured}, \emph{predicted} or \emph{heuristic}, and each is followed by the number that tests it.

\paragraph{Setting.}
A crop $I$ shows an object with rotation $q\in SO(3)$ (allocentric where \cref{sec:method} uses it) and right-acting symmetry
group $\mathcal S$ (continuous symmetries discretised), so $q$ and $qs$, $s\in\mathcal S$, produce the same image. The error
is $e(\hat q,q)=\min_{s\in\mathcal S} d(\hat q,qs)$ with $d$ the geodesic angle. The frozen backbone delivers $z=\varphi(I)$;
every head is a function of $z$ only, so its risk is at least the Bayes risk given $z$, which is at least that given $I$
(data processing). A depth-$L$ tree returns the centre of a cell of the depth-$L$ tesseract partition. We write
$\bar m_L$ for the mean own-cell error over Haar rotations and $r_L$ for the largest distance from a cell centre to a point of
its cell. The angle to a fixed quaternion is quasi-convex on each cell (the cell lifts to a convex polytope in $\mathbb R^4$
and $\{v:\angle(v,a)\le\theta\}$ is a convex cone for $\theta<90^\circ$), so $r_L$ is attained at a cell corner and we
computed it exactly over all corners:

\begin{center}\scriptsize\setlength{\tabcolsep}{4pt}\resizebox{\columnwidth}{!}{%
\begin{tabular}{@{}lccccccc@{}}
\toprule
level $L$ & 1 & 2 & 3 & 4 & 5 & 6 & 7\\
\midrule
$\bar m_L$ (mean, $2\!\times\!10^6$ Haar samples) & 39.27 & 19.95 & 10.01 & 5.01 & 2.51 & 1.25 & 0.63\\
$r_L$ (exact, over all cell corners) & 81.79 & 46.83 & 24.43 & 12.36 & 6.20 & 3.10 & 1.55\\
\bottomrule
\end{tabular}}
\end{center}
Per level, $\bar m_L$ halves to within 2\% from level 2 on and $r_L$ to within 1.2\% from level 4 on [measured].

% -----------------------------------------------------------------------------------------------------------------------------
\subsection{Central Result: Selection and Resolution}

\paragraph{Assumption M (two-term model).}
For a fixed object, crop distribution and trained tree, with probability $\pi$ the image is \emph{in mode}: the returned leaf
is the depth-$L$ cell that contains $q\,\Exp(w)$, where $w\in\mathbb R^3$ is a learned error with a law $W$ that does not depend
on $L$. Otherwise the tree makes a \emph{selection error}, whose error law $P_{\mathrm{sel}}$ does not depend on $L$ either.
The model makes one substantive claim, that depth changes only the quantisation of an otherwise fixed in-mode estimate. Its
parameters are fitted, not derived, and the depth sweep tests it (\cref{sec:supp_theory_scaling}).

\begin{theorem}[Selection--resolution decomposition]\label{thm:decomp}
Under Assumption M, let $E_L$ be the mean error and $R_L(\tau)$ the recall at $\tau$ of the depth-$L$ tree, and let $e_\infty=e(q\Exp(w),q)\le\|w\|$
be the in-mode error without quantisation.
\begin{itemize}
\item[(i)] \emph{Decomposition.} $E_L=\pi\,\mathbb E[e_L\mid\text{in mode}]+(1-\pi)\,\mathbb E_{P_{\mathrm{sel}}}[e]$, and only the first term depends on $L$.
\item[(ii)] \emph{Depth acts only on resolution, by at most $\pi r_L$.} For every $L$, $|e_L-e_\infty|\le r_L$ for in-mode images, so
\[
|E_L-E_\infty|\le\pi\,r_L,\qquad |E_{L+1}-E_L|\le\pi\,(r_L+r_{L+1}),
\]
and $R_L(\tau)\ge\pi\,\Pr(e_\infty\le\tau-r_L)+(1-\pi)P_{\mathrm{sel}}(e\le\tau)$
and $R_L(\tau)\le\pi\,\Pr(e_\infty\le\tau+r_L)+(1-\pi)P_{\mathrm{sel}}(e\le\tau)$.
\item[(iii)] \emph{Selection errors are invariant to depth.} No choice of $L$ changes the fraction $1-\pi$ of images or their errors;
in particular, depth cannot remove flips.
\end{itemize}
\end{theorem}

\noindent\emph{Proof.} (i) is the law of total expectation with Assumption M. (ii) In mode, the returned centre $c$ and the
point $q\Exp(w)$ lie in one cell, so $d(c,q\Exp(w))\le r_L$. For every $s$, the triangle inequality gives
$d(c,qs)\le d(q\Exp(w),qs)+r_L$ and $d(q\Exp(w),qs)\le d(c,qs)+r_L$; minimising over $s$ gives $|e_L-e_\infty|\le r_L$. Taking
expectations over in-mode images gives the bounds on $E_L$, and the recall bounds follow from
$\{e_\infty\le\tau-r_L\}\subseteq\{e_L\le\tau\}\subseteq\{e_\infty\le\tau+r_L\}$. (iii) is Assumption M. \hfill$\square$

\paragraph{Remark 1 (saturation) [heuristic].}
Write $e_L\approx\|w+u\|$ with $u$ the quantisation offset in the tangent space. The offset is a deterministic function of
$q\Exp(w)$, not independent noise, but if it behaves like zero-mean noise of mean square $\propto\bar m_L^2$, then for
$\bar m_L\ll\mathbb E\|w\|$ the linear term of $\mathbb E\|w+u\|$ vanishes and the in-mode mean error exceeds $\mathbb E\|w\|$ by
$O(\bar m_L^2/\sigma)$: the gain in the \emph{mean} of a further level falls by about four per level. Depth stops paying for
the mean at $L^\ast$ where $\bar m_L=\mathbb E\|w\|$, $L^\ast=\log_2(78.5^\circ/\mathbb E\|w\|)$, with $\mathbb E\|w\|=1.60\,\sigma$ for
Gaussian $w$ of per-axis deviation $\sigma$. This argument does not apply to recall at thresholds comparable to $r_L$,
which a finer grid moves directly.

\paragraph{Reading.}
The theorem separates what depth can and cannot do. Depth moves in-mode errors by at most $r_L$ ($6.20^\circ$ at $L=5$,
$3.10^\circ$ at $L=6$); it cannot repair a flip, because a flip is an error of the other term. Every other component of the
method acts on $\pi$, which is why the paper is organised around selection. Fitting Assumption M with Gaussian $w$ to the
depth-5 test errors of all 30 T-LESS and 21 YCB-V objects (EXP-131, DINOv3, per object, least squares on the empirical
distribution function below $12^\circ$) gives $\sigma=2.75^\circ$ per axis and $\pi=0.863$ (T-LESS, orbit) and
$\sigma=2.75^\circ$, $\pi=0.868$ (YCB-V, greedy), object-balanced [measured]. The fit reproduces the depth-5 recalls exactly
by construction; it is the depth sweep that tests it. With $\sigma=2.75^\circ$, $\mathbb E\|w\|=4.4^\circ$ and $L^\ast\approx4.2$: for
the mean error and for R@10, our depth-5 trees are past the saturation depth; for R@2 and R@5, whose thresholds are comparable to
$r_5=6.20^\circ$, they are not (\cref{sec:supp_theory_scaling}).

\paragraph{Remark 2 (SPEED+ per level).}
EXP-050 truncated one SPEED+ tree at levels 3, 4 and 5. Its synthetic validation errors were $11.5^\circ$, $7.27^\circ$ and
$5.54^\circ$, against grid means of $10.0^\circ$, $5.0^\circ$ and $2.5^\circ$. In quadrature,
$\sqrt{E_L^2-\bar m_L^2}=5.7^\circ,5.3^\circ,4.9^\circ$. This uses the total mean, flips included, and drifts by 15\% over two levels
while the grid term falls by four, which is roughly what Assumption M states.
On real lightbox images the error barely changed from level 4 to 5 ($27.8^\circ\to26.9^\circ$). By \cref{thm:decomp}(iii) this
is what a selection-dominated error does: the real-domain error of that tree was mostly flips (\cref{sec:spp}).

% -----------------------------------------------------------------------------------------------------------------------------
\subsection{Propositions: Where Each Component Acts}

\begin{proposition}[Selection is decided at the top of the tree]\label{prop:toplevels}
If $e(\hat q,q)>2r_\ell$, the depth-$\ell$ prefix of $\hat q$ differs from that of every equivalent pose $qs$.
In particular, an error above $2r_1=163.6^\circ$ is decided at the root (the chart) or level 1, one above $2r_2=93.7^\circ$
by level 2, and one above $2r_3=48.9^\circ$ by level 3.
\end{proposition}
\noindent\emph{Proof.} Two rotations in one depth-$\ell$ cell are both within $r_\ell$ of its centre, hence within $2r_\ell$ of
each other. \hfill$\square$

\noindent For the test predictions with error above $30^\circ$ (T-LESS orbit, 1619 image--seed pairs; YCB-V greedy, 969), the
first level at which the predicted path leaves the path of the nearest equivalent ground truth is the chart in 60.5\% and 76.2\%,
level 1 in 27.7\% and 16.6\%, level 2 in 10.2\% and 7.0\%, level 3 in 1.7\% and 0.2\%, and never a deeper level [measured]. A
geometric null, the ground truth rotated by the same angle about a random axis, gives almost the same distribution (58.2, 28.5,
10.9, 2.3\% and 77.1, 18.8, 3.6, 0.5\%). The early divergence is therefore a consequence of the size of the errors, not evidence
about the tree. What the proposition does establish is that for 98--99.8\% of these errors the wrong decision is one among at
most $4\cdot8^2=256$ cells. A tree trains each such decision on every image, and a flat classifier over the same 131{,}072
leaves has to make it among all leaves at once. We suggest this as the reason for EXP-127, where a flat classifier on the
identical leaves reaches $5.34^\circ$ against the tree's $3.87^\circ$ under the same recipe [heuristic]: a level-$\ell$ cell receives
on average $n/(4\cdot8^{\ell})$ training images, about $8^{5-\ell}$ times as many as a flat leaf.

\begin{proposition}[Multi-path targets]\label{prop:multipath}
Let $T(q)$ be the soft target of \cref{sec:tree} at $q$, $T_{\mathrm{multi}}(q)=\frac{1}{|\mathcal S|}\sum_sT(qs)$ the
multi-path target of \cref{eq:symloss}, and $T_{\mathrm{can}}(q)=T(\kappa(q))$ the canonical target, with $\kappa(q)$ the
member of $\{qs\}$ closest to the identity. For a roll $\rho=\Exp(-\theta d)$ acting on the left (\cref{eq:roll}):
\begin{itemize}
\item[(a)] $T_{\mathrm{multi}}(qs')=T_{\mathrm{multi}}(q)$ for every $s'\in\mathcal S$ (invariance), and $T_{\mathrm{multi}}(\rho q)=\rho_\ast T_{\mathrm{multi}}(q)$
(equivariance), exactly, because $(\rho q)s=\rho(qs)$.
\item[(b)] $T_{\mathrm{can}}$ is invariant but not equivariant: $\kappa(\rho q)\ne\rho\,\kappa(q)$ whenever the roll moves the closest-to-identity member
to another member, which for a nontrivial $\mathcal S$ happens on a set of positive measure. For Haar-random poses and a $90^\circ$ roll, $\kappa(\rho q)\ne\rho\,\kappa(q)$ for 49\% of poses with a 2-fold and 91\% with a 4-fold symmetry axis ($2\times10^5$ samples) [measured].
\item[(c)] The multi-path loss is the single-label loss averaged over the annotated member, $\mathcal L_{\mathrm{sym}}(q)=\mathbb E_{s\sim U(\mathcal S)}\mathcal L_{\mathrm{GT}}(qs)$.
If the annotated member is uniform on $\mathcal S$ given the image, both losses have the same expectation and the same minimiser, and the
multi-path loss is its Rao--Blackwellised, lower-variance version. If it is not, as with a fixed object frame and a non-uniform
pose distribution, the single-label minimiser is the posterior of the dataset's annotation convention, which no image determines.
\end{itemize}
\end{proposition}
\noindent\emph{Proof.} (a) Associativity of the quaternion product, and invariance of the average over a group under
right multiplication by a group element. (b) $\kappa$ depends on the distance of each member to the identity, which the left
roll changes; for a nontrivial $\mathcal S$ the region where two members are equally close to the identity is a hypersurface,
and a roll by $\theta\ne0$ carries a set of positive measure across it. (c) Linearity of expectation; the variance statement is
the Rao--Blackwell inequality applied to the gradient. \hfill$\square$

\noindent Orbit voting maps a candidate from view $k$ back with $\rho^{-1}$ and scores it in every view
(\cref{eq:orbit}). With a canonical tree, the mapped-back candidate is in general a non-canonical member, to which the
other views assign almost no mass; by (a), a multi-path tree assigns it the same mass as the member it came from. This is the
mechanism we propose for the measured interaction of EXP-106 [measured, mechanism heuristic]: canonical labels help a single view
but break orbit voting (SYMSOL-T orbit $9.73^\circ$ against $2.07^\circ$ for multi-path targets; T-LESS validation orbit
$10.71^\circ$ against $9.25^\circ$ for single labels).

\begin{proposition}[Orbit voting]\label{prop:orbit}
Let $I_k$ be the crop rolled by $\theta_k$, $z_k=\varphi(I_k)$, and $p_k$ the tree's distribution in view $k$, mapped back to the
original frame.
\begin{itemize}
\item[(a)] \emph{Information.} Let $z_k$ collect the features of all branches in use. Each $I_k$ is a function of $I$, but $z_k$ is
not a function of $z_0$, because the frozen backbones and the pooling are not roll-equivariant. Hence the Bayes risk given
$(z_0,\dots,z_{K-1})$ can be lower than that given $z_0$: for a fixed set of backbones, orbit voting is the only component of the
method that changes what they deliver. A single-view student
(\cref{sec:distill}) sees only $z_0$, so it can recover the part of the orbit gain that is estimation error, but not the part that
is information, even with unlimited data.
\item[(b)] \emph{Pooling bound.} For the normalised geometric pool $\bar p\propto\prod_kp_k^{1/K}$ and every pose $y$,
$-\log\bar p(y)\le\frac1K\sum_k-\log p_k(y)$, with equality if and only if all $p_k$ are equal. The pool is never worse in log-loss
than the \emph{average} view, but it can be worse than view 0 when the rolled views are worse than view 0.
\item[(c)] \emph{Prior.} If $p_k(c)\propto\ell_k(c)\,p(c)$ with a common prior, the score $s(c)=\sum_k\log p_k(c)$ counts the prior $K$ times.
\end{itemize}
\end{proposition}
\noindent\emph{Proof.} (a) As stated; the Bayes risk can only decrease when conditioning on more variables. (b) By the
inequality of arithmetic and geometric means, $Z=\sum_y\prod_kp_k(y)^{1/K}\le\sum_y\frac1K\sum_kp_k(y)=1$, so
$-\log\bar p(y)=\frac1K\sum_k-\log p_k(y)+\log Z\le\frac1K\sum_k-\log p_k(y)$. (c) Algebra. \hfill$\square$

\noindent The statements concern the normalised pool and the score, not the arg-max over 32 candidates that we report, and
not the weighted vote. They give a measurable condition for when voting can hurt: the rolled views must be worse than the
unrolled one. We measured each view's best beam candidate, mapped back, on the EXP-131 depth-5 trees [measured]:

\begin{center}\scriptsize\setlength{\tabcolsep}{3pt}
\begin{tabular}{@{}llcccc@{}}
\toprule
& & \multicolumn{2}{c}{view 0} & \multicolumn{2}{c}{views 1--3}\\
\cmidrule(lr){3-4}\cmidrule(lr){5-6}
Dataset & split & mean & R@10 & mean & R@10\\
\midrule
T-LESS (30) & val & $8.67^\circ$ & 79.5 & $8.74$--$8.86^\circ$ & 79.2--79.4\\
            & test & $13.56^\circ$ & 77.3 & $14.64$--$14.78^\circ$ & 72.9--74.3\\
YCB-V (21)  & val & $8.00^\circ$ & 84.0 & $11.79$--$12.18^\circ$ & 62.7--64.2\\
            & test & $9.06^\circ$ & 83.6 & $11.45$--$12.39^\circ$ & 64.2--67.9\\
\bottomrule
\end{tabular}
\end{center}
On T-LESS validation the four views are equally good and voting lowers the error ($8.54^\circ\to6.84^\circ$). On YCB-V the
rolled views lose about 20 points of R@10, and voting raises the validation error ($7.58^\circ\to8.71^\circ$), as on the paper's
six objects (EXP-122). On the YCB-V test set the mean nonetheless falls ($11.05^\circ\to9.54^\circ$), because voting removes some
gross errors (above $90^\circ$: 2.9\% $\to$ 1.9\%) while it loses fine accuracy (R@5 49.8\% $\to$ 45.1\%). Why rolled YCB-V
views are worse we have not tested; household objects filmed upright give the tree an orientation prior that rolling breaks,
and (c) shows that the score then counts a misplaced prior four times [heuristic].

\begin{proposition}[Allocentric conversion]\label{prop:allo}
Let the egocentric estimate be $R_{\mathrm{ray}}(t)\hat A$, with $\hat A$ the allocentric estimate and $R_{\mathrm{ray}}$ the minimal
rotation taking the optical axis to the ray of $t$. Replacing $t$ by $\hat t$ changes the egocentric error by at most
$d\big(R_{\mathrm{ray}}(\hat t),R_{\mathrm{ray}}(t)\big)$, for every image and every symmetry set.
\end{proposition}
\noindent\emph{Proof.} The geodesic distance is bi-invariant, so $d(R_{\mathrm{ray}}(\hat t)\hat A,R_{\mathrm{ray}}(t)\hat A)=d(R_{\mathrm{ray}}(\hat t),R_{\mathrm{ray}}(t))$;
apply the triangle inequality to the error against each equivalent ground truth and minimise over $\mathcal S$. \hfill$\square$

\noindent With the predicted translations of EXP-131, $d(R_{\mathrm{ray}}(\hat t),R_{\mathrm{ray}}(t))$ has mean $0.07^\circ$
(95th percentile $0.16^\circ$, maximum $0.56^\circ$) on T-LESS and $0.10^\circ$ ($0.24^\circ$, $0.58^\circ$) on YCB-V (seed 0) [measured].
The bound explains why EXP-129 found that predicted translation changes the mean rotation error by at most $0.01^\circ$: the
largest possible per-image change is a fraction of a degree, and its sign varies between images.

% -----------------------------------------------------------------------------------------------------------------------------
\subsection{Remarks}

\paragraph{Remark 3 (frozen features bound what the heads can reach).}
Every head conditions on $z=\varphi(I)$, so the gap to fine-tuned methods has two parts: the information that $\varphi$ discards,
which only training the backbone recovers, and the heads' estimation error, which our components reduce. On the same frozen
DINOv3 features the heads close 30--58\% of the gap between a frozen baseline and LoRA fine-tuning (\cref{sec:ft}); in this
reading that is the estimation part. Orbit voting is the exception of \cref{prop:orbit}(a): it changes $z$ without training $\varphi$, as does adding a second
backbone [heuristic].

\paragraph{Remark 4 (product of experts).}
The per-level PoE adds the log-probabilities of the two branches. It is the posterior given both features if the features are
conditionally independent given the path prefix, and it then counts the prior twice, by the algebra of \cref{prop:orbit}(c).
Fusion can only help selection if the branches' errors fall on different modes. We have not measured how often they do, and this
remark does not explain why fusion helps on SPEED+ and T-LESS but not on SYMSOL-T and YCB-V.

\paragraph{Remark 5 (distillation and orbit voting).}
By \cref{prop:orbit}(a), a single-view student cannot retain the information part of the orbit gain, which is consistent with
the measured retention of 13--43\% on SPEED+. A student trained towards the orbit teacher also makes the four views agree more,
which leaves the vote less to correct; this is consistent with the slightly worse SYMSOL-T orbit system after distillation
[heuristic; view agreement before and after distillation is not yet measured].

% -----------------------------------------------------------------------------------------------------------------------------
\subsection{Scaling Laws}
\label{sec:supp_theory_scaling}

\paragraph{Depth.}
Under Assumption M with Gaussian $w$, the depth-$L$ error follows from the depth-5 error and the fitted $(\pi,\sigma)$, exactly
within the model:
\[
E_L=E_5+\pi\big(\mathbb E[e_L]-\mathbb E[e_5]\big)_{\text{in mode}},\qquad
R_L(\tau)=R_5(\tau)+\pi\big(\Pr(e_L\le\tau)-\Pr(e_5\le\tau)\big)_{\text{in mode}},
\]
with the in-mode terms simulated with the exact tesseract encoder, per object, at that object's test poses and symmetry set.
For the log-likelihood at the ground truth (EXP-131 protocol) we additionally assume that the tree's in-mode density is calibrated, a Gaussian of the
same $\sigma$ around its estimate; the in-mode change is then $\pi\,\big(\mathbb E\log P(\text{GT leaf})-\mathbb E\log\mathrm{vol}(\text{GT leaf})\big)\big|^{L}_{5}$.
The leaf volume falls by 8 per level, so a tree whose in-mode mass stays in one leaf gains $\ln8=2.08$ nats per level; once
cells are smaller than $\sigma$, the mass spreads over several leaves and the gain vanishes.

\paragraph{A test that needs no new trees [measured].}
Greedy decoding is sequential, so the depth-$\ell$ prefix of a depth-5 greedy leaf is exactly the tree's greedy depth-$\ell$
decision. Fitting $(\pi,\sigma)$ on the depth-5 greedy test errors and predicting the truncated trees with the formula above gives,
object-balanced, at depth 4: T-LESS mean $12.81^\circ$ against $12.62^\circ$ measured, R@5 25.0 against 28.1, R@10 72.1 against 73.5;
YCB-V $12.40^\circ$ against $12.38^\circ$, 27.9 against 28.7, 77.6 against 78.1. At depth 3 the means are $16.27^\circ$ against
$15.79^\circ$ and $16.00^\circ$ against $16.30^\circ$. The model over-predicts R@2 at coarse depths (YCB-V depth 4: 4.0 against 2.4).
This tests the shape of the model and the simulation, not the assumption that separately trained deeper trees keep the
same $W$; only the depth sweep tests that.

\paragraph{Prediction for EXP-132 [predicted].}
Fitted on the EXP-131 depth-5 test errors of the same trees that EXP-132 re-evaluates at depth 5, and stated before any depth-6 or
depth-7 tree was evaluated (object-balanced, DINOv3; test split):

\begin{center}\scriptsize\setlength{\tabcolsep}{3pt}
\begin{tabular}{@{}llccccc@{}}
\toprule
Dataset (mode) & $L$ & mean & R@2 & R@5 & R@10 & $\Delta$LL vs.\ $L5$\\
\midrule
T-LESS, 30 (orbit)  & 5 (measured) & $10.90^\circ$ & 6.1 & 48.7 & 83.8 & 0\\
                    & 6 & $10.47^\circ$ & 8.5 & 57.1 & 84.7 & $+0.27$\\
                    & 7 & $10.34^\circ$ & 9.4 & 59.6 & 84.9 & $+0.42$\\
\midrule
YCB-V, 21 (greedy)  & 5 (measured) & $11.05^\circ$ & 7.9 & 49.8 & 83.6 & 0\\
                    & 6 & $10.62^\circ$ & 10.9 & 57.8 & 84.4 & $+0.31$\\
                    & 7 & $10.51^\circ$ & 11.9 & 59.7 & 84.6 & $+0.44$\\
\bottomrule
\end{tabular}
\end{center}
The law predicts most of the gain at depth 6, at the tight thresholds (R@2, R@5), almost none at R@10 and in the share of gross
errors, and a log-likelihood gain far below $\ln8$ per level. If the depth-6 and depth-7 trees learn a smaller in-mode error than
the depth-5 trees, the measured gains will exceed these; if their deeper levels learn worse, as they see fewer images per
cell, the gains will fall short. The log-likelihood prediction is the least secure: the depth-6 and depth-7 soft targets are
$1.25^\circ$ and $0.63^\circ$ wide, narrower than $\sigma$, so the deeper trees may be over-confident, and their log-likelihood can
then fall short of the calibrated prediction, or fall, while their rotation error improves. Either outcome is reported against this table. SYMSOL-T is not predicted: the depth-5
evaluation of its trees did not save per-image errors.

\paragraph{Views [heuristic].}
If the per-view log-likelihood margins between the true mode and a flip had mean $\mu>0$, variance $v$ and pairwise correlation
$\rho$, the probability that the flip wins the sum over $K$ views would fall as $\exp\!\big(-K_{\mathrm{eff}}\mu^2/2v\big)$ with
$K_{\mathrm{eff}}=K/(1+(K-1)\rho)$, which saturates at $1/\rho$. Rolled views of one crop are strongly correlated, so few views
should capture most of the gain. Evaluating $K=2,4,8$ on existing trees would test this; we have not done so.

% -----------------------------------------------------------------------------------------------------------------------------
\subsection{Scope and Observations That Do Not Fit}
\begin{itemize}
\item \emph{Assumption M is a model.} It is fitted ($\sigma,\pi$ per object) and tested only by the depth sweep. Gaussian $w$ is
used only for Remark 1 and for the prediction; \cref{thm:decomp} does not need it. The log-likelihood prediction also needs calibration.
\item \emph{Proposition 1 is a necessary condition.} That wrong predictions leave the correct path at the chart in most cases is
measured, not implied.
\item \emph{Tree against flat.} The sample-count argument after \cref{prop:toplevels} is a heuristic. A data-scaling experiment
(training on subsets) would test it; we have not run one.
\item \emph{Orbit voting.} \Cref{prop:orbit}(b) concerns the normalised pool, not the arg-max we report, and the cause of the
worse rolled views on YCB-V is not tested. On T-LESS test the rolled views are also somewhat worse than view 0 (R@10 about 4
points lower), yet voting helps; the condition is necessary for harm, not sufficient.
\item \emph{Unexplained.} Why geometry fusion helps only on two benchmarks (Remark 4), and why the information and estimation parts of the
orbit gain have the sizes they have (Remark 5).
\end{itemize}
